\chapter{Erectile Dysfunction (Impotence)}

\section*{Definition}
The persistent or recurrent inability to achieve or maintain an erection sufficient for satisfactory sexual activity, when it causes personal or relationship distress. Duration and severity are assessed in context rather than by a single cutoff.

\section*{When to See a Doctor}
See your doctor if:
\begin{itemize}
  \item Sudden onset of ED after pelvic trauma or surgery, ED with loss of morning erections (possible hypogonadism), or ED causing significant relationship distress
\end{itemize}

\section*{How It Develops}
Erection requires intact neurological, vascular, endocrine, and psychological systems. Nitric oxide released from cavernous nerve terminals and endothelium activates guanylate cyclase, increasing cGMP, causing smooth muscle relaxation and arterial dilation. ED results from disruption at any level: endothelial dysfunction (atherosclerosis, diabetes) impairs NO production; venous leak prevents trapping of blood; neurological damage disrupts signaling; low testosterone reduces libido and erectile function.

\section*{Causes}
\begin{itemize}
  \item Vascular disease (atherosclerosis, hypertension, diabetes mellitus)
  \item Neurological disorders (spinal cord injury, multiple sclerosis, Parkinson, stroke)
  \item Endocrine causes (hypogonadism, hyperprolactinemia, thyroid disorders)
  \item Psychological factors (anxiety, depression, performance anxiety, stress)
  \item Medications (antidepressants SSRIs, antihypertensives beta-blockers, diuretics, antipsychotics)
  \item Lifestyle factors (smoking, alcohol, obesity, sedentary lifestyle, recreational drugs)
  \item Pelvic surgery or trauma (prostatectomy, cystectomy, pelvic radiation)
  \item Peyronie disease (penile curvature with plaques)
  \item Renal failure or liver cirrhosis
  \item Aging-related decline in testosterone and vascular function
\end{itemize}

\section*{Related Symptoms}
\begin{itemize}
  \item Low libido
  \item Premature ejaculation
  \item Difficulty achieving orgasm
  \item Penile curvature
  \item Urinary symptoms
\end{itemize}

\section*{Medical Evaluation}
\begin{itemize}
  \item Your doctor will take a detailed sexual, medical, and psychosocial history.
  \item Testing commonly includes a morning total testosterone level and may include glucose or lipids; prolactin, thyroid testing, or repeat testosterone testing is targeted to the history and examination.
  \item A physical examination includes genital examination and assessment of secondary sexual characteristics.
  \item Nocturnal penile tumescence testing or penile Doppler ultrasound may be used to differentiate organic from psychogenic causes.
  \item Referral to a urologist or endocrinologist is arranged for complex or refractory cases.
\end{itemize}

\section*{Treatment Options}
\begin{itemize}
  \item Oral phosphodiesterase-5 inhibitors are commonly first-line pharmacotherapy when appropriate; they must not be combined with nitrates and require review of cardiovascular status and other medicines.
  \item Testosterone replacement therapy for confirmed hypogonadism.
  \item Psychosexual counseling for psychogenic ED or relationship issues.
  \item Vacuum erection devices or penile injections (alprostadil) as second-line options.
  \item Penile implant surgery for severe, treatment-resistant cases.
\end{itemize}


\section*{Self-Care and Home Management}
\begin{itemize}
  \item Engage in regular aerobic exercise to improve vascular health and blood flow.
  \item Maintain a healthy weight and follow a heart-healthy diet (Mediterranean diet).
  \item Limit alcohol consumption and stop smoking to improve erectile function.
  \item Reduce stress through mindfulness, meditation, or counseling.
  \item Communicate openly with your partner about concerns and expectations.
\end{itemize}

\section*{References}
\begin{thebibliography}{99}
\bibitem{erectile_dysfunction:shamloul2013}
Shamloul R, Ghanem H. Erectile dysfunction. \textit{Lancet}. 2013;381(9861):153--165.
\bibitem{erectile_dysfunction:nehra2015}
Nehra A, Jackson G, Miner M, et al. The Princeton III consensus for management of erectile dysfunction and cardiovascular disease. \textit{Mayo Clin Proc}. 2012;87(8):766--778.
\bibitem{erectile_dysfunction:montorsi2010}
Montorsi F, Adaikan G, Becher E, et al. Summary of the recommendations on erectile dysfunction. \textit{J Sex Med}. 2010;7(11):3467--3478.
\bibitem{erectile_dysfunction:hiatz2014}
Hatzichristou D, Rosen RC, Br\"{o}derick G, et al. Clinical evaluation and management strategy for sexual dysfunction in men and women. \textit{J Sex Med}. 2004;1(1):49--57.
\bibitem{erectile_dysfunction:mcintosh2000}
McIntyre RS, Mancini D, Eisfeld BS, et al. Erectile dysfunction. \textit{N Engl J Med}. 2000;342(24):1802--1813. [Published correction in \textit{N Engl J Med}. 2000;343(3):224].
\bibitem{erectile_dysfunction:akerlof2001}
Akerl\"{o}f E, Sj\"{o}strand I. Erectile dysfunction. In: \textit{International Encyclopedia of Urology}. Wiley; 2015.
\end{thebibliography}
